% Flavor catalog per-process draft.
% process_id: L001
% family: charged_lepton
% owner_agent_id: PKA-L001
% writer_agent_id: null
% checker_agent_id: null
% opus_signoff_id: null
% Canonical metadata sidecar: L001.yaml
% Status is controlled by the YAML sidecar status_history, not this comment.

\section{L001: \texorpdfstring{\(\mu\to e\gamma\)}{mu -> e gamma} Dipole LFV}

\subsection*{Process}
\textbf{Standard notation:} \(\mathcal{B}(\mu^+\to e^+\gamma)\).
This entry covers the charged-lepton-flavor-violating radiative muon decay and
the dipole-style bound on the off-diagonal \(e\mu\) lepton-flavor spurion used
by the current RS lepton-sector scan.  It is a charged-lepton process, not part
of the quark \(\Delta F=2\) lane.

\subsection*{PDG-or-equivalent value}
The current primary-experiment limit is the MEG II 2025 result recorded in
\texttt{L001.yaml:primary\_current\_limit},
\[
  \mathcal{B}(\mu^+\to e^+\gamma) < 1.5\times 10^{-13}\quad (90\%\ {\rm C.L.}),
\]
from the 2021--2022 MEG II data set, with quoted search sensitivity
\(2.2\times10^{-13}\).  The local snapshot is
\texttt{flavor\_catalog/references/L001/megii2025\_arxiv2504\_15711.txt}.
For audit context, \texttt{L001.yaml:pdg\_2025\_listing} records the PDG 2025
muon-listing value
\[
  \mathcal{B}(\mu^+\to e^+\gamma) < 3.1\times10^{-13}\quad (90\%\ {\rm C.L.}),
\]
which is the MEG+MEG II 2021 combined limit and appears to lag the later MEG II
2025 publication.

\subsection*{Relevance to RS / anarchic flavor}
In the lepton extension of the warped model, \(\mu\to e\gamma\) probes
brane-localized or loop-induced dipoles controlled by the off-diagonal
combination \((\bar Y_N\bar Y_N^\dagger)_{12}\).  The repo implements the
Perez--Randall NDA form
\(\mathcal{B}\simeq 4\times10^{-8}\,|(\bar Y_N\bar Y_N^\dagger)_{12}|^2
(3\,{\rm TeV}/M_{\rm KK})^4\), converts an experimental branching-ratio limit
to \(C=\sqrt{\mathcal{B}_{\rm lim}/(4\times10^{-8})}\), and uses the updated
MEG II limit in the scan default.  Those numerical inputs are the
\texttt{L001.yaml:paper\_era\_reference} and \texttt{L001.yaml:repo\_default}
entries: the live default gives \(C\simeq1.936\times10^{-3}\), compared with
the paper-era \(C=0.02\).

\subsection*{Post-2008 developments}
Relative to the contemporaneous Perez--Randall lepton paper, the experimental
story changed substantially.
Perez--Randall used the MEGA-era
\(\mathcal{B}(\mu\to e\gamma)<1.2\times10^{-11}\) input
(\texttt{L001.yaml:paper\_era\_reference}).  The full MEG 2009--2013 data set
then gave \(4.2\times10^{-13}\), MEG II's 2021-only result gave
\(7.5\times10^{-13}\), and the MEG+MEG II combination gave
\(3.1\times10^{-13}\) (\texttt{L001.yaml:prior\_experimental\_limits}).  The
MEG II 2021--2022 analysis now sets the \(1.5\times10^{-13}\) primary value.

\subsection*{Constraint validity / model dependence}
The experimental bound is clean: a signal would be charged-lepton flavor
violation with negligible Standard Model background.  Its translation into this
repo is model-dependent, because it relies on the chosen RS lepton spurion, NDA
prefactor, and \(M_{\rm KK}\) convention rather than a full loop calculation.
The catalog should therefore classify it as a robust experimental limit but a
lepton-extension-only, dipole/NDA model interpretation.

\subsection*{Code coverage in this repo}
\textbf{YES.}  The sidecar \texttt{code\_coverage.evidence} block lists the
implementation in \texttt{flavorConstraints/muToEGamma.py}: the module
documents the \(\mu\to e\gamma\) dipole constraint, records the paper-era
\(C=0.02\) coefficient, defines
\(\texttt{PREFAC\_BR}=4.0\times10^{-8}\), converts branching-ratio limits to
\(C\), and exposes \(\texttt{check\_mu\_to\_e\_gamma}\).  The scan default is the
MEG II 2025 value:
\texttt{scanParams/scan.py}:32--33 sets
\(\texttt{BR\_LIMIT\_MEGII\_2025}=1.5\times10^{-13}\), and line 264 makes it the
default \(\texttt{ScanConfig.br\_limit}\).  Tests at
\texttt{tests/test\_mu\_to\_e\_gamma.py}:40 and
\texttt{tests/test\_scan.py}:50 exercise the updated limit.

\subsection*{Implementation difficulty}
\textbf{LOW} for cataloging and present code coverage, because the observable
is already implemented and the scan default matches the primary 2025 result.  A
future production refresh is \textbf{MEDIUM}: the code should audit the
paper-era NDA normalization, \(M_{\rm KK}\) convention, and whether the final
catalog should expose both paper-reproduction and MEG II 2025 modes.

\subsection*{Key references}
Process-local source keys before bibliography consolidation:
\texttt{MEGII2025\_MuEGamma}, \texttt{MEGII2024\_MuEGammaFirst},
\texttt{MEG2016\_MuEGammaFull}, \texttt{PDG2025\_MuonListing}, and
\texttt{PerezRandall2008\_WarpedNeutrinos}.
